\documentclass[11pt]{article}

\usepackage[margin=1in]{geometry}
\usepackage[T1]{fontenc}
\usepackage{amsmath,amssymb,amsthm,mathtools}
\usepackage{microtype}
\usepackage{booktabs,tabularx}
\usepackage{enumitem}
\usepackage{xcolor}
\usepackage[colorlinks=true,linkcolor=blue!55!black,
  citecolor=blue!55!black,urlcolor=blue!55!black]{hyperref}


\setlist[itemize]{topsep=2pt,itemsep=2pt}

\numberwithin{equation}{section}
\newtheorem{theorem}{Theorem}[section]
\newtheorem{lemma}[theorem]{Lemma}
\newtheorem{proposition}[theorem]{Proposition}
\newtheorem{corollary}[theorem]{Corollary}
\newtheorem{conjecture}[theorem]{Conjecture}
\theoremstyle{definition}
\newtheorem{example}[theorem]{Example}
\newtheorem{remark}[theorem]{Remark}

\newcommand{\cO}{\mathcal O}
\newcommand{\cZ}{\mathcal Z}
\newcommand{\Rcnt}{\mathsf R}
\newcommand{\Mcnt}{\mathsf M}
\newcommand{\avgdeg}{\overline d}

\title{A detailed version of the recent proof of the Erd\H{o}s--S\'os theorem}
\author{Louis DeBiasio}
\date{September 6, 2026}

\hypersetup{pdfauthor={Louis DeBiasio}, pdftitle={A marked-position proof of the Erdos-Sos theorem}}

\begin{document}
\maketitle

\begin{abstract}
We give a detailed exposition of the marked-position proof of the
Erd\H{o}s--S\'os theorem: every graph of average degree greater than $t-2$
contains every tree on $t$ vertices.  The argument counts permutations
with an adjacency-marked prefix and bounds the number of states that
fail to support a prescribed rooted tree.  A shortest-prefix rotation
joins branches, and a prefix reversal adds a new leaf.  We give both
injections and their inverse maps explicitly, beginning with the path case.
\end{abstract}

\section*{Acknowledgments}
ChatGPT~6 Astra was used to develop, write, and revise this exposition.
The author's role was minimal, consisting of asking questions while
trying to understand the recent Erd\H{o}s--S\'os proof.
The original Erd\H{o}s--S\'os proof and its benchmark provenance are
recorded in \cite{AdamczewskiBloom2026,Erdos548Repository}.

\section{Introduction}

This paper expounds the marked-position argument in the Lean proof of
Erd\H{o}s problem~548 produced by GPT-6 Astra in the FrontierMath
Erd\H{o}s benchmark \cite{AdamczewskiBloom2026,Erdos548Repository}.
The proof presented here is due to that work; the purpose of this exposition
is to present its finite counting argument in conventional mathematical form.
The conjecture itself is due to Paul Erd\H{o}s and Vera T.~S\'os.
It appears in Erd\H{o}s's 1964 paper \cite[p.~30]{Erdos1964}, where he
explicitly attributes it jointly to S\'os and himself; the paper is authored
by Erd\H{o}s alone.

All graphs in this note are finite and simple, and a copy need not be
induced.  We use $t$ for the number of vertices of the target tree.

\begin{theorem}[Erd\H{o}s--S\'os]\label{thm:ES}
Let $T$ be a tree on $t\geq 2$ vertices and let $G$ be a graph on $n$
vertices.  If $G$ contains no copy of $T$, then
\[
        e(G)\leq \frac{(t-2)n}{2}.
\]
Equivalently, every graph of average degree greater than $t-2$ contains a
copy of every tree on $t$ vertices.
\end{theorem}

The bound is sharp.  A disjoint union of copies of $K_{t-1}$ has average
degree $t-2$ and has no connected subgraph on $t$ vertices.  The star gives a
second obstruction: a $(t-2)$-regular graph contains no copy of
$K_{1,t-1}$.

If minimum degree were greater than $t-2$, the result would be immediate:
root $T$ anywhere and embed one new leaf at a time.  Average degree is much
less cooperative.  Deleting low-degree vertices guarantees only a nonempty
subgraph of minimum degree at least roughly half the average degree, losing
the factor of two that matters here.  Much of the earlier literature is
devoted to recovering that loss by understanding the structure of the host
graph or of the target tree.  For example, Besomi, Pavez-Sign\'e, and
Stein \cite{BesomiPavezSigneStein2021} proved the conjecture for bounded-degree
trees in large dense host graphs.  Pokrovskiy \cite{Pokrovskiy2024} proved a
structural hyperstability theorem for graphs excluding a bounded-degree tree,
and applied it to establish the conjecture for sufficiently large
bounded-degree trees without a density restriction on the host.
These results provide background; neither is used as a black box below.

The new proof never constructs a single embedding greedily.  Instead, it
runs a rooted embedding test simultaneously through every ordering of
$V(G)$ and amortizes the possible failures exactly.

\section{A warm-up: the Erd\H{o}s--Gallai path theorem}
\label{sec:path-warmup}

For paths, the counting argument uses only prefix reversal.  We give
this case first to illustrate the method without the additional operation
needed to join branches.

\begin{theorem}[Erd\H{o}s--Gallai \cite{ErdosGallai1959}]
\label{thm:EG-path}
Let $t\geq2$.  If an $n$-vertex graph $G$ contains no path on $t$
vertices, then
\[
        e(G)\leq\frac{(t-2)n}{2}.
\]
\end{theorem}

\begin{proof}
The case $n=0$ is immediate, so assume $n\geq1$.
For $j\geq2$, let $A_j$ count the pairs $(\pi,i)$ such that
$\pi=(v_1,\ldots,v_n)$ is an ordering of $V(G)$, $2\leq i\leq n$,
$v_1v_i\in E(G)$, and the prefix $\{v_1,\ldots,v_i\}$ contains a
path on $j$ vertices with endpoint $v_1$.  Every pair satisfying the
adjacency condition supplies a two-vertex path.  Thus
\[
        A_2=\sum_{x\in V(G)}d_G(x)(n-1)!=2e(G)(n-1)!.
\]
We claim that
\begin{equation}
        A_j\leq A_{j+1}+n!\qquad(j\geq2).
        \label{eq:path-warmup-step}
\end{equation}

For each ordering that contributes to $A_j$, discard its first counted
pair $(\pi,i_0)$.  There are at most $n!$ discarded pairs.  For every
remaining pair $(\pi,i)$, we have $i>i_0$, so the earlier prefix
contains a $j$-vertex path with endpoint $v_1$ that avoids $v_i$.
The edge $v_iv_1$ extends this to a $(j+1)$-vertex path with endpoint
$v_i$.  Reverse the prefix through position $i$:
\[
 (\pi,i)\longmapsto
 \bigl((v_i,v_{i-1},\ldots,v_1,v_{i+1},\ldots,v_n),i\bigr).
\]
The image is counted by $A_{j+1}$: its first vertex is $v_i$, its
marking edge is still $v_iv_1$, and its prefix has the same vertex set.
This map is injective, since the index $i$ is retained and reversing
the prefix again recovers the original ordering.  This proves
\eqref{eq:path-warmup-step}.

Iterating from $j=2$ to $j=t-1$ gives
\[
        2e(G)(n-1)!=A_2\leq A_t+(t-2)n!.
\]
Since $G$ contains no path on $t$ vertices, $A_t=0$.  Dividing by
$2(n-1)!$ gives the result.
\end{proof}

The same reversal will extend a rooted tree by a new leaf in
Lemma~\ref{lem:leaf}.  For a general tree, we also need an injection
that combines branches at a common root.  We next introduce the notation
for these two operations.

\section{Set-up}

For the state-counting argument, assume $n\geq1$; the empty host
satisfies Theorem~\ref{thm:ES} immediately.  Fix an ordering
\[
        \pi=(v_1,v_2,\ldots,v_n)
\]
of $V(G)$.  For $1\leq i\leq n$, the \emph{position $i$} may be displayed by
inserting a divider after $v_i$:
\begin{equation}
 (v_1,v_2,\ldots,v_i\mid v_{i+1},\ldots,v_n).
 \label{eq:position-display}
\end{equation}
If $i=n$, the right-hand side of the divider is empty.  The vertices to
the left of the bar form the \emph{prefix at position $i$}.

For $2\leq i\leq n$, call this position \emph{marked} when
$v_1v_i\in E(G)$.  In other words, the first and last vertices of the prefix
are adjacent.  The vertical bar carries no graph structure; it merely records
which prefix is under consideration.  Let
\[
 \cO_G=\{(\pi,i):\pi\text{ is an ordering of }V(G),\ 2\leq i\leq n,
                         \ v_1v_i\in E(G)\}.
\]
Write $\Mcnt(G)=|\cO_G|$.

\begin{proposition}\label{prop:M}
For every graph $G$ on $n\geq1$ vertices,
\[
        \Mcnt(G)=2e(G)(n-1)!.
\]
Consequently $\Mcnt(G)/n!=\avgdeg(G)$, the average degree of $G$.
\end{proposition}

\begin{proof}
If the first vertex is $x$, each of the $(n-1)!$ orderings of the remaining
vertices has exactly $d_G(x)$ marked positions.  Sum over $x$ and use the
handshaking lemma.
\end{proof}

Now let $(H,h)$ be a rooted graph.  A marked state $(\pi,i)$ \emph{supports}
$(H,h)$ if the set
\[
        \{v_1,v_2,\ldots,v_i\}
\]
contains a copy of $H$ in which $h$ is mapped to $v_1$.  Define
\[
 \cO_G(H,h)=\{(\pi,i)\in\cO_G:(\pi,i)\text{ supports }(H,h)\}.
\]
The host graph will always be clear, so from now on we abbreviate
$\cO_G(H,h)$ to $\cO(H,h)$.

\begin{example}[Supportive and failed states]\label{ex:marked}
Let $G$ be the graph with
\[
        V(G)=\{a,b,c,d,e\}
\]
and edge set
\[
        ac,\quad ae,\quad ad,\quad eb
\]
and consider
\[
        \pi=(a,c,e,b,d).
\]
The positions after $c,e,d$ are marked, because these vertices are adjacent
to the first vertex $a$; the position after $b$ is not marked.  Thus the
four positions behave as follows:
\[
\begin{array}{ccl}
 (a,c\mid e,b,d)     &:& \text{marked, but its prefix has only one edge},\\
 (a,c,e\mid b,d)     &:& \text{marked, and contains the path $c-a-e$},\\
 (a,c,e,b\mid d)     &:& \text{unmarked, and contains the path $a-e-b$},\\
 (a,c,e,b,d\mid)     &:& \text{marked, and contains the path $a-e-b$}.
\end{array}
\]
Let $H$ be a three-vertex path rooted at an endpoint.  The position after $e$
does not support $(H,h)$: its prefix contains a copy of $H$, but in that
copy the first vertex $a$ is the middle of the path rather than the image of
the root.  The position after $b$ contains the correctly rooted path
$a-e-b$, but is not marked.  Only the final position is both marked and
supportive.
Thus this ordering contributes three states to $\Mcnt(G)$ and one state to
$|\cO(H,h)|$.
\end{example}

\begin{remark}[Why the marked edge need not be used]
The mark and the rooted copy play deliberately different roles.  In the
final state of Example~\ref{ex:marked}, the edge $ad$ certifies that the
position is marked, while the rooted path is $a-e-b$; the edge $ad$ is not
part of that path.  A state should therefore be viewed as carrying two
pieces of information:
\begin{itemize}
\item the prefix has already made a rooted copy of the target available;
\item the last vertex of the prefix supplies one neighbour-token at the
      root.
\end{itemize}
The second item is what connects the permutation count to density.  For a
fixed ordering with first vertex $x$, there is exactly one marked position for
each neighbour of $x$, so the number of marked positions is $d_G(x)$.  If
every position $2,\ldots,n$ were counted instead, every ordering would
contribute $n-1$ states, independently of the edges of $G$, and the count
would contain no average-degree information.

The unused marking edge is also useful in the injections.  In the leaf move,
the vertex at a later marked position can be attached to an earlier rooted
copy as the new leaf.  In the branch rotation, the same edge remains the
anchor that makes the new position marked.  Thus the proof does not assert
that every marking edge belongs to the copy being counted; it uses marked
edges as transferable opportunities and counts those occurring after a
rooted copy has become available in $|\cO(H,h)|$.
\end{remark}

The heart of the proof is the following stronger, rooted statement.

\begin{theorem}[Master inequality]\label{thm:master}
Let $(T,r)$ be any rooted tree on $t\geq2$ vertices.  For every host graph
$G$ on $n\geq1$ vertices,
\begin{equation}
        \Mcnt(G)\leq|\cO(T,r)|+(t-2)n!.
        \tag{MI}\label{eq:MI}
\end{equation}
Equivalently,
\[
        \avgdeg(G)\leq\frac{|\cO(T,r)|}{n!}+t-2.
\]
\end{theorem}

If $G$ is $T$-free, then $\cO(T,r)=\varnothing$.  Proposition~\ref{prop:M} and
\eqref{eq:MI} immediately give Theorem~\ref{thm:ES}.  Everything therefore
reduces to proving the master inequality.

The term $|\cO(T,r)|/n!$ is the expected number of marked positions, in a
uniformly random ordering, whose prefixes contain a copy of $T$ rooted at
the first vertex.  The marking edge need not belong to that copy.  Therefore
\[
        \frac{\Mcnt(G)-|\cO(T,r)|}{n!}
\]
is the expected number of marked positions whose prefixes do not contain
such a rooted copy.  The master inequality says that this expectation is at
most $t-2$.

\section{The two injections}\label{sec:injections}

\subsection{Joining two branches by rotating a word}

\begin{lemma}[Branch gluing]\label{lem:branch}
Let $G$ be an $n$-vertex graph.  Suppose $T_1$ and $T_2$ are rooted
subtrees of a tree $T$, all rooted at $r$, such that
\[
 T_1\cup T_2=T,\qquad V(T_1)\cap V(T_2)=\{r\},
\]
where the union is taken as subgraphs of $T$.  Then
\[
 |\cO(T_1,r)|+|\cO(T_2,r)|
 \leq \Mcnt(G)+n!+|\cO(T,r)|.
\]
\end{lemma}

\begin{proof}
For the rotation it is convenient to adjoin the root-only state to every
ordering:
\[
        \cZ=\{(\pi,1):\pi\text{ is an ordering of }V(G)\}.
\]
Thus $|\cZ|=n!$.  A root-only state is a bookkeeping device, not a marked
state.

We define an injection
\begin{equation}
 \Phi:\cO(T_1,r)\setminus\cO(T,r)
 \longrightarrow
 \bigl(\cO_G\setminus\cO(T_2,r)\bigr)\cup\cZ.
 \label{eq:branch-injection}
\end{equation}

Take $(\pi,i)$ in the domain.  Because it supports $T_1$, there is a first
marked position $a\leq i$ whose prefix supports $T_1$.
Divide the entries of $\pi$ following $v_1$ into three consecutive blocks:
\[
 \pi=(v_1,A,B,C),
 \qquad
 A=(v_2,\ldots,v_a),
 \quad B=(v_{a+1},\ldots,v_i),
 \quad C=(v_{i+1},\ldots,v_n).
\]
Here the commas denote concatenation of blocks; they are not dividers.  The
original state has index $i$, so its divider appears after the block $B$.

Let $b=1+|B|=i-a+1$, and form a new ordering by interchanging the
consecutive blocks $A$ and $B$:
\[
        \rho=(v_1,B,A,C).
\]
Define
\begin{equation}
        \Phi(\pi,i)=(\rho,b).
 \label{eq:rotation}
\end{equation}
Thus the image uses position $b$, whose divider lies immediately after $B$;
its prefix is
$v_1$ followed by the entries of $B$.  If $B$ is empty, then $b=1$ and the
image lies in $\cZ$.  Otherwise the vertex at the new position is the last
vertex of $B$, namely $v_i$.  Hence position $b$ is marked, because the
source position $i$ was marked and $v_1v_i\in E(G)$.

The new prefix, consisting of $v_1$ and the entries of $B$, cannot support
$T_2$.  Indeed, $v_1$ together with the entries of $A$ contains a rooted
copy of $T_1$.  A rooted copy of $T_2$ using only $v_1$ and the entries of
$B$ would be disjoint from it except at $v_1$; the two copies could
therefore be joined at their common root to give a copy of $T$ in the
original prefix, which consists of $v_1,A,B$.  This contradicts the fact
that $(\pi,i)$ does not support $T$.  Thus $\Phi(\pi,i)$ belongs to the
stated codomain.

It remains to show that $\Phi$ is injective.  Given an image $(\rho,b)$,
the block $B$ is already known: it consists of the $b-1$ entries following
the first vertex.  Starting immediately to the right of the divider, let $A'$
be the shortest nonempty block such that
\begin{enumerate}[label=(\roman*),leftmargin=2.3em]
\item the last vertex of $A'$ is adjacent to the first vertex of $\rho$;
      and
\item $A'$ together with the first vertex supports rooted $T_1$.
\end{enumerate}
The original block $A$ has these properties.  No proper initial subblock of
$A$ can have them, because it would have produced an earlier marked position
supporting $T_1$ in the original ordering, contrary to the choice of $a$.
Hence $A'=A$.  We can therefore recover
\[
        \pi=(v_1,A,B,C),\qquad i=1+|A|+|B|,
\]
so $\Phi$ is injective.

Its codomain has size $\Mcnt(G)-|\cO(T_2,r)|+n!$.  Since every state
supporting $T$ also supports $T_1$, the domain has size
$|\cO(T_1,r)|-|\cO(T,r)|$.  Therefore
\[
 |\cO(T_1,r)|-|\cO(T,r)|
 \leq \Mcnt(G)-|\cO(T_2,r)|+n!,
\]
which is the claimed inequality.
\end{proof}

\begin{example}[A branch rotation]\label{ex:rotation}
Let $T_1$ be a single edge rooted at one endpoint, and let $T_2$ be a
two-edge path rooted at an endpoint.  Their union at the root is a
four-vertex path rooted at an internal vertex.  Consider the fragment
\[
        \pi=(o,a,x,y,\ldots),
        \qquad A=(a),\quad B=(x,y),
\]
where the state uses the divider after $y$.  Suppose
that, among the displayed vertices, the only edges are $oa$ and $oy$.  The block $A$ together with $o$ already supports $T_1$, while the
prefix through $y$ does not support $T$: neither $x$ nor $y$ continues to a
second edge.

Here the state uses position $i=4$, since four vertices lie to the left of
the divider, and the first supporting position is $a=2$.  Thus
$b=1+|B|=3$.  The map $\Phi$ interchanges $A$ and $B$, producing
\[
        \rho=(o,x,y,a,\ldots)
\]
with marked position $3$, again immediately after $y$.  It is marked
because $oy$ is an edge, but its prefix $o,x,y$ does not support $T_2$.
To invert the move, scan to the right of the divider; the shortest block ending
at a neighbour of $o$ and supporting rooted $T_1$ together with $o$ is
precisely $A=(a)$.
\end{example}

\subsection{Moving the root across a new leaf}

\begin{lemma}[Leaf move]\label{lem:leaf}
Let $G$ be an $n$-vertex graph and let $(S,p)$ be a rooted tree.  Form $T$
by adding a new leaf $\ell$ adjacent to $p$, and root the larger tree at
$\ell$.  Then
\[
        |\cO(S,p)|\leq|\cO(T,\ell)|+n!.
\]
\end{lemma}

\begin{proof}
Let
\[
        \cZ=\{(\pi,1):\pi\text{ is an ordering of }V(G)\}
\]
be the set of root-only states.  We define an injection
\[
        \Psi:\cO(S,p)\longrightarrow \cO(T,\ell)\cup\cZ.
\]

Fix an ordering $\pi$.  If $\pi$ contributes any states to $\cO(S,p)$,
let $i_0$ be the smallest index for which $(\pi,i_0)\in\cO(S,p)$, and set
\[
        \Psi(\pi,i_0)=(\pi,1)\in\cZ.
\]
Thus the first supporting state belonging to each ordering is sent to that
ordering's unique root-only state.

Now take any other state $(\pi,i)\in\cO(S,p)$.  Since it is not the first
such state for $\pi$, we have $i_0<i$.  In particular, position $i_0$ is an
earlier marked position whose prefix supports a copy of $S$ rooted at $v_1$.
The endpoint $v_i$ of the current prefix is adjacent to $v_1$, and it is
not used by the earlier
copy of $S$.  Thus adjoining $v_i$ to $v_1$ creates a copy of $T$ in which
the new root $\ell$ is represented by $v_i$.

Reverse the prefix through $v_i$, leaving the rest of the ordering fixed:
\begin{equation}
 \rho=(v_i,v_{i-1},\ldots,v_2,v_1,
               v_{i+1},\ldots,v_n),
 \qquad
 \Psi(\pi,i)=(\rho,i).
\label{eq:reversal}
\end{equation}
The image uses the same position $i$.  Its new first vertex is $v_i$ and its
last vertex is $v_1$, so it is marked.  Its prefix has exactly the same
vertex set as the original prefix; it therefore contains the earlier copy
of $S$ together with the edge $v_iv_1$.  Hence $\Psi(\pi,i)$ supports
$(T,\ell)$.

For each fixed $i$, the map reverses the first $i$ entries and is therefore
an involution.  Since the position index is part of the state, images with
different $i$ cannot collide.  These images lie in $\cO(T,\ell)$, whereas
the exceptional images lie in $\cZ$; and the latter are distinct because
they retain their orderings.  Thus $\Psi$ is injective.  Since
$|\cZ|=n!$, we obtain
\[
        |\cO(S,p)|\leq|\cO(T,\ell)|+n!,
\]
as claimed.
\end{proof}

\begin{example}[The leaf involution]\label{ex:leaf}
Let $S$ be an edge rooted at $p$, and let $T$ be the three-vertex path
$\ell-p-q$, rooted at the new leaf $\ell$.  Consider
\[
        \pi=(b,u,x,c,d)
\]
with edges $bu$ and $bc$.  The position after $u$ is the first marked
position supporting $S$, so it is sent to the root-only state associated
with this ordering.  Now consider the later state
$(\pi,4)$, whose position is displayed as
\[
        (b,u,x,c\mid d).
\]
The index is $4$ because four vertices lie to the left of the divider (in
particular, $c=v_4$).  This state also supports $S$.  Using $b-u$ as the
earlier copy,
attach $c$ to $b$ and reverse the prefix through $c$.  Thus
\[
        \rho=(c,x,u,b,d),
        \qquad
        \Psi(\pi,4)=(\rho,4),
\]
and the image position is displayed as
\[
        (c,x,u,b\mid d).
\]
Its prefix contains the path $c-b-u$, rooted at its first vertex $c$, and
the position is marked because its first and last vertices, $c$ and $b$, are
adjacent.
\end{example}

\section{The main inductive proof}

We now prove Theorem~\ref{thm:master}.  The two injections account for the constant $t-2$ in the induction.

\begin{proof}[Proof of Theorem~\ref{thm:master}]
We induct on $t=|V(T)|$, uniformly in the host graph and in the chosen root.

If $t=2$, then $T$ is an edge.  Every marked position itself exhibits a rooted
copy of that edge, so $\Mcnt(G)\leq|\cO(T,r)|$.  This is \eqref{eq:MI} with
no error term.

Suppose $t\geq3$ and first assume that $r$ is a leaf.  Let $p$ be its
neighbour, and root $S=T-r$ at $p$.  By induction and Lemma~\ref{lem:leaf},
\[
\begin{aligned}
 \Mcnt(G)\leq |\cO(S,p)|+(t-3)n!\leq |\cO(T,r)|+(t-2)n!.
\end{aligned}
\]

It remains to treat the case $d_T(r)\geq2$.  Choose a neighbour $s$ of $r$
and another neighbour $z\neq s$.  Delete the edge $rs$, and let $A$ be the
component containing $r$.  Define
\[
        T_1=T[A],\qquad
        T_2=T[(V(T)\setminus A)\cup\{r\}],
\]
both rooted at $r$.  Both trees have at least two vertices and are smaller
than $T$.  They satisfy the hypotheses of Lemma~\ref{lem:branch}.  If
$t_j=|V(T_j)|$, then the common root is counted twice, so
\[
        t_1+t_2=t+1.
\]
Induction gives
\[
\begin{aligned}
 \Mcnt(G)&\leq |\cO(T_1,r)|+(t_1-2)n!,\\
 \Mcnt(G)&\leq |\cO(T_2,r)|+(t_2-2)n!.
\end{aligned}
\]
After addition,
\[
 2\Mcnt(G)\leq
 |\cO(T_1,r)|+|\cO(T_2,r)|+(t-3)n!.
\]
Lemma~\ref{lem:branch} now yields
\[
        2\Mcnt(G)\leq\Mcnt(G)+|\cO(T,r)|+(t-2)n!,
\]
and subtracting $\Mcnt(G)$ proves \eqref{eq:MI}.
\end{proof}

Combining this with Proposition~\ref{prop:M} completes the proof of
Theorem~\ref{thm:ES}.  In the customary notation, if $T$ has $k+1$
vertices and
\[
        e(G)>\frac{k-1}{2}n,
\]
then $G$ contains a copy of $T$; equivalently, since $e(G)$ is integral, it
suffices that
\[
        e(G)\geq \left\lfloor\frac{(k-1)n}{2}\right\rfloor+1.
\]


\begin{thebibliography}{99}

\bibitem{AdamczewskiBloom2026}
T.~Adamczewski and T.~F. Bloom,
\emph{FrontierMath Erd\H{o}s}, Epoch AI, 2026.
See in particular Appendix~B.4 and the accompanying FrontierMath Erd\H{o}s
benchmark materials, \url{https://epoch.ai/files/frontiermath-erdos.pdf}.

\bibitem{BesomiPavezSigneStein2021}
G.~Besomi, M.~Pavez-Sign\'e, and M.~Stein,
\emph{On the Erd\H{o}s--S\'os conjecture for trees with bounded degree},
Combin. Probab. Comput. \textbf{30} (2021), no.~5, 741--761.
\href{https://arxiv.org/abs/1906.10219}{arXiv:1906.10219}.

\bibitem{Erdos1964}
P.~Erd\H{o}s,
\emph{Extremal problems in graph theory},
in \emph{Theory of Graphs and its Applications}
(Proc. Sympos. Smolenice, 1963),
Publishing House of the Czechoslovak Academy of Sciences, Prague, 1964,
29--36.  The joint conjecture with V.~T.~S\'os is stated on p.~30.
\href{https://users.renyi.hu/~p_erdos/1964-06.pdf}{Original paper}.

\bibitem{ErdosGallai1959}
P.~Erd\H{o}s and T.~Gallai,
\emph{On maximal paths and circuits of graphs},
Acta Math. Acad. Sci. Hungar. \textbf{10} (1959), 337--356.
\href{https://users.renyi.hu/~p_erdos/1959-10.pdf}{Original paper}.

\bibitem{Erdos548Repository}
T.~Adamczewski,
\emph{Erd\H{o}s problem \#548: Lean proof and provenance}, 2026.
\url{https://github.com/tadamcz/erdos548}.

\bibitem{Pokrovskiy2024}
A.~Pokrovskiy,
\emph{Hyperstability in the Erd\H{o}s--S\'os Conjecture},
preprint, 2024.
\href{https://arxiv.org/abs/2409.15191}{arXiv:2409.15191}.

\end{thebibliography}
\end{document}
